% =====================================================================
%  前置部分：译校说明 + 编者前言（まえがき，原书 p.i–ii）
% =====================================================================
\chapter*{译校说明}
\addcontentsline{toc}{chapter}{译校说明}
\markboth{译校说明}{译校说明}

\section*{一、底本}

本书译自深谷賢治 編《ミラー対称性入門》（日本評論社，2009 年 9 月 15 日 第 1 版第 1 刷，ISBN 978-4-535-78632-5），其底本是『数学セミナー』2006 年 10 月号至 2007 年 8 月号的连载，成书时作了若干增补与修改（其中第 3 章为新写）。

这是一本由九位作者分头执笔的入门书，不是深谷賢治一人的专著。全书十二章，执笔者依次为：江口徹（第 1 章）、深谷賢治（第 2、8、10 章）、高橋篤史（第 3、5、11 章）、並河良典（第 4 章）、細野忍（第 6 章）、今野宏（第 7 章）、Martin Guest 与乙藤隆史（合著第 9 章）、中島啓（第 12 章）。各章自成一格，文风、记号与体例并不统一；中译本如实保留这种差异，不强行拉平。

原书正文 163 页，另有索引（p.164–165）、初出一覧（p.166）、著者紹介（p.167）。初出一覧与编者紹介已照译，收入书末；原书索引按日文五十音排列，对中文读者检索价值有限，中译本未逐条转录，改以书末“附录 B：术语表（日汉对照）”代替——该表按主题分为十一组，组内大致按日文五十音顺序排列，条目数超过原书索引。

\section*{二、体例}

\begin{enumerate}[label=\arabic*.,leftmargin=2.2em,itemsep=2pt]
  \item \textbf{页码对照}：每页正文开头以 \texttt{\% ===== PDF p.N (原书 p.M) =====} 标注该页对应的 PDF 页码与原书页码；图注末尾标“（原书 p.N）”。原书页码 = PDF 页码 − 8（PDF 第 9 页即原书 p.1）。
  \item \textbf{公式与陈述编号}：原书公式按每章重置的（1）（2）……编号，图、表同样每章重置。定义、定理、引理、命题、推论、例、注等陈述的编号也随原书，各章体例不一（有的章编号、有的章不编号）。中译本一律照原编号排，不重编、不补号。（第 4 章 p.57 的两幅插图原书本无图号，为便于称引，中译本补编为图 1、图 2，并在批注中说明，这是全书唯一的补号例外。）
  \item \textbf{节号}：原书各章节号体例亦不统一（用罗马数字或阿拉伯数字）。中译本按原书逐章处理，不改动。
  \item \textbf{图}：原书插图均自扫描页按 400--600 dpi 裁出，图号与图题照原书，图内的日文标注保留原样（不重绘、不改字），中文图题由译者另加并注明原书页码；不裁入原书的日文图题行，以免与中文图题重复。
  \item \textbf{译校批注}：以源码注释 \texttt{\% [Tutor Note: ...]} 的形式给出，不进入正文排版，全部汇集于书末“附录 A：译校批注”，按章排列并注明原书页码。批注只写三类内容：易错预警（符号、指标、约定）、作者省略论证的点拨、背景直觉；另有少数几条说明图版与体例（原书无图号、原书表无题等）。原书疑似笔误一律照原样排，只在批注中指出并给出正确形式。
  \item \textbf{参考文献}：各章末尾的参考文献照录原书，日文与英文条目均不译（第 3、10、11 章原书末页即接下一章，没有参考文献页，译本同样不设）。
  \item \textbf{作者读音标记}：原书各章末尾印有作者姓名的假名读音（如“［ふかやけんじ］”），译注本省略。
  \item \textbf{本章要点框与译者评注框}：版式上另有“本章要点”“译者评注”两种方框（均非原书内容），本书未使用。
\end{enumerate}

\section*{三、术语原则}

术语以书末“附录 B：术语表（日汉对照）”为准。原则是：中文学界有稳定、无歧义通行译名者用中文；无稳定译名、或硬译反而妨碍与文献对照者保留原文（例如 flop、crepant、nef、holonomic D-模、Gromov--Witten 不变量一类专名与缩写）。人名一律不译，只按日文原书读音或英文原拼写列出。日文与中文用法有别者按中文数学惯例处理，例如：日文“多様体”译“流形”，“層”译“层”，“圏”译“范畴”，“複体”译“复形”，“特異点”译“奇点”，“因子”在代数几何语境译“除子”、在一般语境译“因子”。

\section*{四、声明}

本书为学术研读用译本，不作商业用途，不对外发行。译文与批注中的错误由译者负责，恳请读者指正。原书的一切著作权归各位作者与日本評論社所有。

\clearpage

% =====================================================================
\chapter*{前言}
\addcontentsline{toc}{chapter}{前言（原书 p.i–ii）}
\markboth{前言}{前言}

本书是把『数学セミナー』的连载（2006 年 10 月号到 2007 年 8 月号）汇集起来、并作了若干增补修改而成的书，主题是镜像对称。

镜像对称本来是理论物理学家发现的现象，但本书的执笔者几乎都是数学家，内容也是数学。请把它当作一本“作为数学的镜像对称”的书来读。

镜像对称最初是作为关于 Calabi--Yau 流形的现象被发现的（现在它的视野已经更宽了）。镜像对称中出现的 Calabi--Yau 流形，典型的是（复）3 维的图形。1 维复流形（黎曼面）按曲率分为三类：正曲率（球面）、零曲率（环面）、负曲率（有两个以上洞的曲面）。曲率为零的情形，也就是环面，其高维推广就是 Calabi--Yau 流形。

把黎曼面以及它上面的种种数学放在一起看，就知道环面（即椭圆曲线）的情形格外丰富。从环面上的函数论（椭圆函数论）、作为复流形的环面的分类（椭圆曲线模空间的研究），直到数论性质的研究，话题非常多。这些话题在 18、19 世纪被认为是数学的最前沿与中心，当时许多数学家都参与其中（数论方面等等，即使现在其重要性也没有改变）。

当（复）维数变成 2 或 3 时，要寻找能产生像椭圆曲线那样丰富数学的空间：2 维的情形是 K3 曲面，3 维的情形则以 3 维 Calabi--Yau 流形为好的候选——我想现在人们相当广泛地相信这一点。对于这些流形，模空间的理论，以及可看作椭圆函数论高维化的种种理论，正在发展（Calabi--Yau 流形的数论性质的研究也许还太难，进展不多）。镜像对称是构成这场研究之中心的现象，现代数学的许多领域都与它相关。

作为数学，或许应当看成一个“猜想”；但这个“猜想”的内容现在并没有被确定下来，构成它一部分的有趣问题接连出现，值得研究的事情不断增加。

越挖越有趣的东西堆成一座宝山——这就是笔者等人对镜像对称的印象。总之，即使不考虑作为其发现背景的物理应用，可研究的对象也取之不尽。

关于镜像对称的书，专业的、数学的、物理立场的，都在逐渐增多。另一方面，镜像对称涉及的数学很多；而且这些数学大多是 20 世纪高度发展起来的、对初学者来说学起来很吃力的东西。等把必要的预备知识全部学完再来做镜像对称的数学，恐怕只会在学预备知识的过程中老去，或者学不下去而放弃。

\enlargethispage{9\baselineskip}
因此我想，像本书这种不太专业的解说书，也许有用。话虽如此，本书是数学书；即使没有严格的证明，所写的内容也没有偷工减料，是意识到数学书应有的正确性来写的。如果书中的各条记述能成为通向各种领域的入口、引导读者走下去，那就太好了。

\medskip
\begin{flushright}
2009 年 8 月\\
深谷 賢治
\end{flushright}
